\section{互补存储机制与原生逻辑矩阵}

采用 \sid{MRAM-06} 的 2T2MTJ 互补存储和 IBMD 本地数字化机制：每个权重 bit 的两个 MTJ 保存 $W$ 与 $\overline W$，P低阻表示0，AP高阻表示1。原40nm宏的每个IBMD含局部latch，其输出是 $\mathrm{OUT}=\mathrm{IN}\land W$。本文保留该器件／bitcell机制，以共同28nm数字外围资源组织有限读码保持与归约；原540MHz@1.2V、最佳能效@0.65V和不同电压的读错误结果分别保留。

\sid{MRAM-06} p.5给出64个bank，每bank有256行、4个权重bit列。两个bank合并形成一个INT8输出，因此本例原生逻辑矩阵为
\[
W[N,K],\quad K=256,\quad N=32,\quad b_S=b_R=1\ \mathrm{Byte}.
\]
每向量输入 $B_S=256$ Byte，有效容量 $KN=8192$ Byte。物理容量与地址为
\[
64\times256\times4=65536\ \text{互补bitcell},\qquad
2\times65536=131072\ \text{MTJ},
\]
\[
W_{o,i}[b]\longmapsto (\mathrm{bank}=2o+\lfloor b/4\rfloor,\;\mathrm{row}=i,\;\mathrm{column}=b\bmod4).
\]
物理MTJ状态共16384Byte，但互补编码不使逻辑payload翻倍。没有增强信号的权重复制。INT8输出所需容器为 $16+\lceil\log_2 256\rceil=24$ bit，交付32个结果。

采用原生bank容量不等于直接采用原宏全部外围。原宏为256项全并行归约；本参考每轮只活动32项、16个输出，数字精度保持，活动位宽和外围资源按共同政策声明。选择八个行组和两个输出组，每tile为 $32\times16\times8=4096$ 个IBMD节点。原bank的4列是四个权重位，不能被误认为能从一个列端口独立读出32行权重。原256项归约树不作为同时工作的服务资源；本参考以16条32项归约替代其服务通路，若保留原树则旁路隔离，不计其额外吞吐。

\subsection{证据的分工}
\sid{MRAM-06}提供原生结构、局部数字化、先双P后单AP的写入和逐row读验机制。\sid{MRAM-03}为28nm 1T1MTJ存储宏，支持完整20/30ns写access及读时序量级；\sid{MRAM-04}为28nm 2T2MTJ差分存储读，补充1.3ns的读量级。\sid{MRAM-01}是电阻串联求和／TDC机制，其11.1MHz与左右顺序写不进入主结果。共同数字预算不将这些来源拼成同芯片实测读写宏。

\clearpage
\section{读码保持的实际通路与输入服务}

\subsection{从输入触发的逻辑值到有限权重寄存}
IBMD原有latch随IN变化：IN为0时输出被拉低，即使W为1也不能把此输出当作已缓存W。原文p.5的read mode令IN由0变1，此时每个本地OUT等于W；原宏再经bank加法树读出选中行结果。这里明确增加\textbf{加法树之前的4096路OUT抽头、tile选择和隔离}，从65536个本地节点中选择一份4096bit tile。八行组×两输出组对应每捕获通道16选1，不能通过读取bank求和值重建独立权重。

以IN固定为1感测所选节点，完成一次物理读后，用一$T_D$捕获到新增4096-bit单bank权重寄存器。随后将阵列输出隔离，由\textbf{4096个独立数字AND}把保持的W与真实输入位相乘，接16条32项归约／累加通道。AND和归约共同使用每轮一$T_D$的数字功能预算；这需要相应逻辑和布线，不声称原IBMD已经免费提供该数字复用路径。

当前tile全部八个输入位用完后才替换。16个24-bit部分和寄存器跨行组保持；输入寄存2048bit、完整输出寄存768bit。原有65536个IBMD latch属于bitcell机制，不作为输入无关的全矩阵shadow使用。更新和终验占用阵列选择时，streaming停止；没有读写同时到峰值的假定。

\subsection{物理读槽、位宽和完整计数}
选择 $t_m=3/5/10$ ns：\sid{MRAM-04}的1.3ns宏级读使用0.3V位线偏置；\sid{MRAM-03}的2.8ns读对应1.2V、25$^\circ$C及有限错误条件。它们为本地读提供量级，不直接保证4096活动节点与新抽头的时序。$t_m$覆盖物理读初始化／预充、差分判决和选通布线；新增权重捕获和隔离交接另付一拍。下一tile感测之前的预充在其读槽内完成。

八个行组、两个输出组构成16个tile；每tile八输入位，故
\[
N_{\rm rd}=N_{\rm cap}=16,\qquad N_D=16\times8=128,
\]
\[
\boxed{\Delta_S=16t_m+16T_D+128T_D+2T_D.}
\]
末两拍为输入接收和完整输出提交。典型 $t_m=T_D=5$ ns，得 $80+80+640+10=810$ ns，$\rho=316.05$ MB/s。数字AND／归约约占79\%，物理读只有约10\%。输入payload较大、输出维度较小以及位复用共同决定该点，不能解释为MTJ材料读速同比提高。

原生256项累加使用24-bit容器；32项局部结果保留符号并继续跨八组累加。无ADC、TDC或量化截断。服务顺序执行，$\Delta_S$为完整向量间隔；快速bitcell延迟本身不会消除多输入位和有限数字通道的占用。

\clearpage
\section{完整两相写入、供流与绝对状态验证}

\subsection{一个更新组的真实资源}
一次对齐更新为一个输入行的八个INT8权重，即64互补对、128MTJ，$B_R=8$ Byte。目标逻辑64bit在本地展开128个P/AP目标bit，经128-bit接口一拍装入128-bit目标寄存器。所选16个bank各含四个权重bit列，配置bank／row写使能和未选SL高阻隔离，不能影响同一原生矩阵其余权重。

\sid{MRAM-06} p.2、Fig.2b/d及SI p.4明确先将两个MTJ都置P，再将每对的一侧置AP：第一相SL接写电压、BL/BLB接0；第二相SL接0，目标BL或BLB接写电压。本文每组第一相有128条活动路径，第二相64条，共192个物理状态命令；不是左侧一次、右侧一次，也不是Flash块擦除。

SI中256对的512$I_W$/256$I_W$关系按实际支路缩放为本例128$I_W$/64$I_W$。固定安装128条方向可选写驱动，每条额定200$\mu$A，合计SL供流能力25.6mA；第二相64条活动时相应12.8mA。\sid{MRAM-06} p.14 Extended Data Fig.1(b)的实测准静态I--V在所示约$\pm0.6$V范围内达到约180$\mu$A量级，支持该有限驱动规模。200$\mu$A是本地参考能力条件，不是20/30ns脉冲电流实测或写成功界；实际脉冲偏置、负载与选择必须在此额度内成立。三情景不改变驱动数量或额定能力。

\subsection{两支路终验不能由差分结果替代}
单个差分逻辑输出可能漏掉P/P、AP/AP非法状态。配置64条单端P/AP感测通道及参考、左右支路隔离、128-bit状态寄存器和64个比较器。分别读取左侧64MTJ和右侧64MTJ；每批完整低偏置读$t_m$之后，显式用一$T_D$捕获，再用一$T_D$比较。两个支路均匹配目标、每对一P一AP后，才完成8Byte payload。

这些验证通道和状态寄存器独立列账，不能从原IBMD差分latch免费获得。读槽含验证模式的初始化、选通及感测；捕获和比较分别收费。目标寄存跨有限重试保持，状态寄存逐批更新。原文有错误row重写流程，但未给完整写验周期或固定重试次数。

\subsection{完整服务账本}
$t_W=20/30/30$ ns为单方向完整写槽，以\sid{MRAM-03}的写access量级为依据，包含该相初始化、驱动、切换／自终止及恢复；不重复叠裸脉冲。两相之间另用一$T_D$切换，事务前端／完成共两拍。因此
\[
C=2t_W+T_D+2(t_m+T_D+T_D),\qquad
\boxed{\Delta_R=2T_D+C=2t_W+2t_m+7T_D.}
\]
典型为 $10+60+5+10+10+10=105$ ns。写两相60ns、读回10ns、捕获10ns、比较10ns、跨相5ns和控制10ns均只计一次。自终止支持节能，不意味着外部固定写槽可以直接换成平均MTJ翻转时间。

\clearpage
\section{相容情景、完整装载与结果解释}

% Generated by check_mram.py --emit.
\begin{table}[htbp]\centering\small
\begin{tabular}{@{}lrrrrrr@{}}\toprule
情景 & $t_m/t_W$ (ns) & $\Delta_S$ (ns) & $\Delta_R$ (ns) & $\rho$ (MB/s) & $\tau$ (MB/s) & $\mathrm{RI}^{*}$\\\midrule
乐观 & 3/20 & 340 & 60 & 752.941 & 133.333 & 5.6471\\
典型 & 5/30 & 810 & 105 & 316.049 & 76.190 & 4.1481\\
悲观 & 10/30 & 1620 & 150 & 158.025 & 53.333 & 2.9630\\
\midrule
典型，重写一次 & 5/30 & 810 & 200 & 316.049 & 40.000 & 7.9012\\
\bottomrule\end{tabular}
\caption{条件配对情景，MB/s 为十进制：$B_S=256$ Byte，$B_R=8$ Byte。前三行为单轮完整写验通过的有限预算；末行为恰好一次整组重写且随后成功的占用对照，未纳入前三行情景范围。}
\end{table}


主情景为一次完整两相和绝对终验被接受的有限服务预算；无实测重试分布时不虚构平均成功率或最坏尾延迟。失败事务不产生有效payload，其时间保留。固定参考资源下另列恰好一次整组重写且第二次通过的占用：目标数据已保持，$\Delta_R^{(2)}=2T_D+2C=200$ ns，$\tau=40$ MB/s、$\mathrm{RI}^*=7.9012$，未并入主三情景。

典型 $B_S=256$ Byte、$B_R=8$ Byte，得到
\[
\rho=316.05\ \mathrm{MB/s},\quad\tau=76.19\ \mathrm{MB/s},\quad
\mathrm{RI}^*=\frac{256}{8}\frac{105}{810}=4.1481.
\]
主三情景固定原生64bank、4096活动读节点和抽头、4096-bit保持、AND／归约、写驱动和绝对验证资源。3/5/10ns读、20/30/30ns写槽与共同2/5/10ns数字拍为相容工程配对，并非同芯片测得的统计快慢界。长情景写槽仍30ns，其他完整服务阶段变慢；不凭空增加重试以扩大范围。

\subsection{完整矩阵装载连接 Table II}
原生矩阵为8192Byte，$256\times(32/8)=1024$个对齐事务覆盖全部有效权重，典型
\[
T_R=1024\times105=107520\ \mathrm{ns},\qquad
U^*=\frac{T_R}{\Delta_S}=132.7407=32\mathrm{RI}^*.
\]
一次完整装载后$U$个向量的$Q_S=256U$ Byte、$Q_R=8192$ Byte、$\mathrm{RI}=U/32$。完整装载接口保留每组控制和两相验收，不使用小事务的$\Delta_R/\Delta_S$替代$U^*$。16KiB展示平均成本为$107520\times16384/8192=215040$ns，是两倍原生容量的平均换算，不是一笔原生请求的延迟，也未据此同时缩放$\rho$与$\tau$。

较大的$\rho$来自原生$K=256,N=32$的入口payload、输出工作量及有限保持；更新能力仍由互补两相、完整终验和实际供流决定。不能将该配置与不同$N$的宏按同一计算工作量排名。扩展到算子时按输入共享、分时和更新域重新得到$T_R,\Delta_S$；FFN和Attention仍逐阶段连接Table II。

\clearpage
\section{维护、证据入口与可复算性}

本地服务窗口内无周期刷新或破坏性读恢复，机器数据分别保存相同的维护前／后能力。\sid{MRAM-03}在其短低偏置读取条件下连续读超过$10^6$次未观察扰动，此为读扰动实验，不是写耐久。两相更新使物理MTJ承受不同的切换／半选历史，寿命、保持和读扰动另作可行性判断，不折入瞬时$\tau$。

\sid{MRAM-05} persistent memory模式允许直接改写0/1，ERASE用于命令兼容；其xSPI速率不能代表局部写并行。本文先双P后单AP属于所选互补机制，不推广为所有MRAM的必要步骤。原1T1MTJ写错误地板不当作该互补宏的精确BER，也不进入重试次数公式。

\subsection{原始证据定位}
\begin{table}[htbp]\centering\small
\begin{tabularx}{\textwidth}{@{}p{24mm}X@{}}\toprule
来源 & 本地PDF页序、条件与用途\\\midrule
\sid{MRAM-06} & 2025 \emph{Nature Electronics}，\href{https://doi.org/10.1038/s41928-025-01479-y}{10.1038/s41928-025-01479-y}。pp.2/4 Fig.2互补IBMD、AND和两相；p.5原生64bank与read mode；p.9逐row写验；p.14 Extended Data Fig.1器件I--V；SI p.4两相电流支路。\\
\sid{MRAM-03} & 2019 \emph{JSSC}，\href{https://doi.org/10.1109/JSSC.2018.2872584}{10.1109/JSSC.2018.2872584}。pp.4--7 Figs.9/15/17/19/21完整读写时序、自终止、温度电压及有限错误地板。\\
\sid{MRAM-04} & 2018 \emph{ISSCC}，\href{https://doi.org/10.1109/ISSCC.2018.8310394}{10.1109/ISSCC.2018.8310394}。pp.1--2 Figs.30.3.2--6，28nm互补存储读、三阶段感测和0.3V读位线条件。\\
\sid{MRAM-01} & 2022 \emph{Nature}，\href{https://doi.org/10.1038/s41586-021-04196-6}{10.1038/s41586-021-04196-6}。pp.6--8，电阻求和、TDC和左右顺序写；与主机制分开。\\
\sid{MRAM-05} & Everspin EMxxxLX/B/HR v3.7，2026-07-23，pp.15/41--48；只支持更新语义，不提供本地时序。\\\bottomrule
\end{tabularx}
\caption{原值、采用值和物理桥接详见证据笔记；各源身份和偏置分别保留。}
\end{table}

\subsection{复算与外部判断}
输入文件\path{data/inputs.json}保存原生组织、局部资源、驱动额定能力和完整阶段。原生配置字段分别存逻辑容量和物理MTJ。情景映射字段\path{mapping_interface}由公共API给出完整$T_R$与$U^*$。\path{scripts/check_mram.py}默认只读，\texttt{--emit}更新未取整结果与表格；检查原生bank地址、活动节点抽头、单tile容量、两支路捕获／比较、独立分项算术和源哈希。

外部审阅重点为新4096路抽头／16选1布线的读槽和负载，新增数字AND与32项归约时序，128路写驱动和25.6mA局部SL供给，以及64路绝对验证在所选偏置下的判决裕量。这些条件需电路或硅测收敛；内部计算检查不等于用户验收。
